% The authors' existing Remarks are merged with Part 6c, rather than replaced.
\section{\src{Discussion and conclusion}}\label{sec:discussion}
\subsection{\src{Summary of findings}}\label{sec:findings}

\src{Using 173 IMF programmes approved in 2002–2016, we asked whether approval-day information can predict whether a conditional lending programme is later interrupted. Our answer is that approval-day information predicts interruption only weakly, and none of the three models is reliably better than the others. On the 28 test programmes, balanced accuracy ranges from 0.51 to 0.57 at thresholds of 0.51 for the forest, 0.50 for logistic regression and 0.44 for the network, against 0.500 for the always-interrupted baseline. At a common threshold of 0.50, the random forest (0.567), the logistic regression (0.570) and the neural network (0.572) lie within 0.005 of each other, while one additional correct prediction moves balanced accuracy by about 0.03 to 0.05. Logistic regression has the highest AUC (0.684).}

\src{Using 13 inputs available at programme approval, we fitted a random forest to predict IMF programme interruptions. Country-grouped cross-validation selected 250 trees and three candidate inputs per split. The final model reached a pooled training out-of-fold balanced accuracy of 0.630, but only 0.508 on the 28 later test programmes at the selected threshold of 0.51. Scheduled reviews, extended facilities and access relative to quota had the highest estimated permutation importance, although these rankings do not establish causal effects. Prior actions rank next.}

\subsection{\src{Limitations}}\label{sec:limitations}

\src{\textbf{Parallel and successive programmes.} Each programme is one observation. Blends approved for the same country on the same day share one review schedule and are merged into one programme. Programmes that overlap or follow an earlier programme of the same country are kept as separate observations: 13 programmes (11 of them in the training set) overlap an earlier one, and 43 of the 84 training countries have more than one programme. Their history enters the models through the track-record inputs. A programme cancelled early and replaced by a successor arrangement counts as interrupted, so the label groups failures with some programmes that were superseded. Because programmes of one country are not independent, we tune the models with country-grouped cross-validation; the test set, by contrast, draws 26 of its 28 programmes from countries in the training set, so it measures prediction for later programmes of mostly known borrowers.}

\src{\textbf{Other limitations.} The results are limited by the small test sample, the change in interruption frequency between training and testing, and the use of training validation results for both model and threshold selection. Test classifications are also sensitive to the threshold: moving from 0.50 to 0.51 changes two predictions. Furthermore, scheduled reviews enter the outcome definition, giving their predictive importance a partly mechanical basis. The macroeconomic inputs are staff projections, and a few extreme values, such as Yemen’s 2014 inflation of −17.7\%, may reflect recording issues in MONA. Our inputs do not measure elections, political cohesion or bureaucratic quality (Dreher, 2003; Ivanova et al., 2003). The share of interrupted programmes rises from 44.8\% in the training set to 60.7\% in the test set, and test programmes are larger, longer and less often concessional, so the drop in performance may reflect a change in the population as well as overfitting.}

\subsection{\src{Policy implications}}\label{sec:policy}

\src{Until performance is shown to be stable across further samples, this forest is best regarded as an exploratory model rather than a dependable stand-alone signal for sovereign-risk decisions. At most, such a score could complement staff judgement. The importance of scheduled reviews must not be read as a reason to schedule fewer reviews: part of the link is mechanical, and fewer reviews would lower measured interruption without improving implementation.}

\src{Extended programmes (EFF or ECF) were interrupted less often than others (38\% against 60\% in the training set), mainly because of ECFs (36\%); EFFs (57\%) differ little from SBAs (60\%). Higher completion in extended programmes need not mean better implementation. Long lending relationships can create moral hazard: a borrower that expects waivers, modified targets or lenient reviews has a weaker incentive to implement, and a lender that has supported a country for years may hesitate to interrupt its programme. The IMF’s Independent Evaluation Office examined such prolonged use of IMF resources (Independent Evaluation Office of the International Monetary Fund [IEO], 2002 [p. ?]), and Reinsberg et al. (2022) use the term “dependency trap” in the title of their study. Because our outcome counts completed reviews, it cannot separate compliance from leniency, so high completion in long programmes is weak evidence of success on its own.}

\subsection{\src{Further research}}\label{sec:further-research}

\src{Future work could expand the sample, assess performance over additional chronological splits, and investigate richer approval-day indicators of political commitment and institutional capacity. A natural next step is to examine a different time period. MONA review records start in 2002, so earlier arrangements would need review histories from other IMF sources, while programmes approved after 2016, including those of the COVID-19 period, can be added once they have run their course. A longer sample would allow rolling time-ordered splits, which train on earlier approval years and test on the following ones and would show whether the weak test result reflects the 2014–2016 period, the small test set or a general lack of signal. Preprocessing and model selection should be confined to the training portions of validation splits.}

\src{Further work could add political and institutional inputs (Dreher, 2003; Ivanova et al., 2003). Coding the type of each condition instead of counting conditions would allow a test of Reinsberg et al. (2022), who associate specific types, such as public-sector reform and privatisation, with interruption. Re-estimating the models without the number of scheduled reviews would show how much of the signal is mechanical, and review-level data would allow predicting each review’s completion with the information available at that review.}
